\chapter{Không có đáp án sạch sẽ cho mọi chuyện}
\markboth{Không có đáp án sạch sẽ cho mọi chuyện}{Không có đáp án sạch sẽ cho mọi chuyện}

\section{Nói thật hết có khi tự hại mình không}
\begin{truyen}{Nói thật hết có khi tự hại mình không}{Classroom Talk}
\chuhoa{H}{ọc sinh: ``Người lớn cứ dạy trung thực. Nhưng có lúc nói thật hết chỉ khiến mình thiệt. Vậy nói thật tới đâu là vừa?''}
\textit{Student: ``Adults keep teaching us to be honest. But sometimes telling the whole truth only hurts us. So how honest should a person be?''}\\[4pt]

\teacherpair{Trung thực không có nghĩa là nói hết mọi thứ theo cách làm người khác tổn thương hoặc đẩy mình vào chỗ dại. Nó là không gian dối điều cốt lõi, và biết chọn cách nói có trách nhiệm.}{Honesty does not mean saying everything in a way that hurts others or pushes yourself into a foolish position. It means not lying about what truly matters, and choosing a responsible way to speak.}

\studentpair{Tức là thật nhưng không ngu?}{So basically, be truthful without being foolish?}

\teacherpair{Có thể nói dân dã vậy.}{You could put it that way, yes.}
\end{truyen}

\begin{baihoc}
Sự thật cần lòng can đảm, nhưng cũng cần sự khôn ngoan trong cách nói.
\textit{(Truth requires courage, but also wisdom in how it is spoken.)}
\end{baihoc}

\begin{ghinhoanh}
\textbf{Short English to remember:}

Honesty needs courage.

It also needs wisdom.
\end{ghinhoanh}

\begin{tuhocnhanh}
1. Em có đang lấy trung thực làm cớ để nói bừa không? \textit{(Do you use honesty as an excuse for careless speech?)}

2. Em cần giữ điều thật nào, và nói nó theo cách nào? \textit{(What truth must you keep, and how should you say it?)}
\end{tuhocnhanh}
\ngancach

\section{Làm người tốt mà toàn bị đè thì sống sao}
\begin{truyen}{Làm người tốt mà toàn bị đè thì sống sao}{Classroom Talk}
\chuhoa{H}{ọc sinh: ``Nghe làm người tử tế thì hay. Nhưng nhiều khi hiền quá là bị leo đầu. Vậy sống tốt kiểu gì cho khỏi ngu?''}
\textit{Student: ``Being a good person sounds nice. But when you are too gentle, people step on you. So how do you live well without being naive?''}\\[4pt]

\teacherpair{Làm người tốt không đồng nghĩa với không có ranh giới. Tử tế mà không biết giữ mình thì dễ thành người luôn phải gánh phần thiệt.}{Being a good person does not mean having no boundaries. If you are kind but do not know how to protect yourself, you easily become the one who always carries the loss.}

\studentpair{Vậy tốt mà cứng được không?}{So can you be good and still be firm?}

\teacherpair{Phải vậy mới bền. Mềm lòng nhưng không mềm nguyên tắc.}{You have to be. That is the only durable way. Keep a soft heart, but do not soften your principles.}
\end{truyen}

\begin{baihoc}
Tử tế không phải để ai muốn làm gì mình cũng được.
\textit{(Kindness does not mean allowing others to do anything to you.)}
\end{baihoc}

\begin{ghinhoanh}
\textbf{Short English to remember:}

Kindness needs boundaries.

Soft heart, firm principles.
\end{ghinhoanh}

\begin{tuhocnhanh}
1. Em có ranh giới nào đang quá yếu không? \textit{(Is any personal boundary of yours too weak?)}

2. Em có thể tử tế mà vẫn dứt khoát ở đâu? \textit{(Where can you stay kind while being firm?)}
\end{tuhocnhanh}
\ngancach

\section{Không ai công nhận thì cố tiếp để làm gì}
\begin{truyen}{Không ai công nhận thì cố tiếp để làm gì}{Classroom Talk}
\chuhoa{H}{ọc sinh: ``Nếu mình cố mà chẳng ai thấy, chẳng ai khen, chẳng ai tin, vậy cố tiếp để làm gì?''}
\textit{Student: ``If I keep trying and nobody notices, nobody praises me, and nobody believes in me, then what is the point of continuing?''}\\[4pt]

\teacherpair{Nếu em chỉ sống bằng tiếng vỗ tay, lúc im lặng kéo dài em sẽ chết tinh thần rất nhanh.}{If you live only on applause, then the moment silence lasts too long, your spirit will collapse very quickly.}

\studentpair{Nhưng ai cũng cần được công nhận mà.}{But everyone needs recognition, do they not?}

\teacherpair{Đúng. Nhưng nếu em không tự biết giá trị của việc mình làm, người ngoài im một cái là em gãy.}{Yes. But if you do not know the value of your own work for yourself, one stretch of silence from other people is enough to break you.}
\end{truyen}

\begin{baihoc}
Cần được ghi nhận là bình thường. Sống lệ thuộc hoàn toàn vào ghi nhận là rất nguy hiểm.
\textit{(Wanting recognition is normal. Depending entirely on it is dangerous.)}
\end{baihoc}

\begin{ghinhoanh}
\textbf{Short English to remember:}

Recognition matters.

Inner ground matters more.
\end{ghinhoanh}

\begin{tuhocnhanh}
1. Em có việc đúng nào đang làm chỉ vì lời khen không? \textit{(Is there any good work you do only for praise?)}

2. Nếu không ai vỗ tay, điều gì vẫn đáng để em làm? \textit{(If no one applauds, what is still worth doing?)}
\end{tuhocnhanh}
\ngancach

\section{Chọn an toàn hay chọn điều mình thích khi nhà không dư dả}
\begin{truyen}{Chọn an toàn hay chọn điều mình thích khi nhà không dư dả}{Classroom Talk}
\chuhoa{H}{ọc sinh: ``Nếu nhà em dư dả chắc em dám chọn cái mình thích. Nhưng nhà không có nhiều, vậy em nên chọn an toàn hay chọn cái em muốn?''}
\textit{Student: ``If my family had money, I would probably dare to choose what I love. But since we do not have much, should I choose safety or what I really want?''}\\[4pt]

\teacherpair{Đây là câu khó thật, không có đáp án sạch. Có người chọn an toàn trước để đỡ gánh gia đình. Có người đi đường vòng để sau mới chạm điều mình thích. Quan trọng là em phải trả giá kiểu nào mà em chịu nổi và không đổ lỗi mãi về sau.}{That is a genuinely hard question, and there is no clean answer. Some people choose safety first to reduce the burden on their family. Others take a longer road and only later reach what they love. What matters is deciding which kind of price you can bear without blaming life forever afterward.}

\studentpair{Nghĩa là chọn gì cũng có mất?}{So no matter what I choose, I will lose something?}

\teacherpair{Ừ. Trưởng thành là biết cái giá nào mình chấp nhận trả.}{Yes. Maturity is knowing which cost you are willing to pay.}
\end{truyen}

\begin{baihoc}
Không phải lựa chọn nào cũng cho mình đủ hết. Nhiều khi chỉ là chọn cái mất mình chịu được.
\textit{(Not every choice gives everything. Sometimes you are choosing the loss you can bear.)}
\end{baihoc}

\begin{ghinhoanh}
\textbf{Short English to remember:}

Some choices are costly either way.

Maturity means owning the cost.
\end{ghinhoanh}

\begin{tuhocnhanh}
1. Điều em thật sự không muốn đánh đổi là gì? \textit{(What is the one thing you truly do not want to trade away?)}

2. Nếu phải đi đường vòng, em có sẵn sàng không? \textit{(If you must take a longer path, are you willing?)}
\end{tuhocnhanh}
\ngancach